% ============================================================
% Disponibilidad de datos y código
% ============================================================

\chapter*{Disponibilidad de datos y código}
\addcontentsline{toc}{chapter}{Disponibilidad de datos y código}
\markboth{Disponibilidad de datos y código}{}

Los datos crudos, el código de procesamiento y análisis, las figuras generadas computacionalmente y los archivos fuente de este documento se encuentran disponibles públicamente en los repositorios del proyecto indicados a continuación. Los archivos crudos se conservan sin modificaciones y cada registro mantiene los identificadores de experimento, muestra, posición y tratamiento necesarios para reconstruir su procedencia.

\section*{Herramientas de desarrollo y asistencia}

Las tareas de programación se realizaron principalmente en \textit{Visual Studio Code}. Para modificar y cargar en el ESP32 el código asociado al controlador PID se utilizó \textit{Arduino IDE}. La redacción y edición del documento en \LaTeX{} contó con asistencia de herramientas de inteligencia artificial generativa: \textit{Claude}, de Anthropic, y \textit{Codex}, de OpenAI. Su uso comprendió la generación y revisión de código \LaTeX{}, la resolución de problemas de compilación y el apoyo en tareas de edición y maquetación. Todo el contenido incorporado fue revisado por la autora, quien conserva la responsabilidad sobre las decisiones metodológicas, el análisis, la interpretación de los resultados y la versión final del texto.

\section*{Repositorios}

\begin{description}\raggedright
  \item[Documento y bibliografía.] Fuentes \LaTeX, bibliografía e imágenes:\\
    \small\url{https://github.com/KierPrev/tesis-magnetobiologia-latex}
  \item[Datos y software.] Mediciones originales, datos procesados y programas de análisis y generación de figuras:\\
    \small\url{https://github.com/KierPrev/Magseed}
\end{description}

Las versiones correspondientes a la entrega final se identificarán mediante los \textit{commits} de ambos repositorios. Las instrucciones de instalación, la organización de los datos y los comandos necesarios para reproducir los resultados se documentan en el archivo \texttt{README} del repositorio de software.

\endinput
